\section{Important Note}
This instruction details the curriculum learning framework for training a specific type of PINN 
designed for this research that focuses on the three-body problem. 
For any information regarding the architecture and structure of the neural network, 
please refer to the \texttt{PINN\_architecture\_plan.tex} file.




\begin{comment}
\section{A Multi-Fidelity Curriculum Learning Framework}


\subsection{Rationale for Curriculum + Transfer Learning in Chaotic Systems}

A curriculum learning approach is highly beneficial for training on chaotic three-body dynamics.  Rather than presenting the neural network with the hardest cases upfront (which could confuse the optimizer or lead to bad local minima), we mimic the pedagogy of physics education: start with the basics, then incrementally increase complexity 16 3. Early in training, the network sees only the stable, regular orbits;  the loss landscape on this subset is much smoother and easier to navigate, allowing the model to tune into fundamental gravitational laws (like the $1/r^2$ force law, energy and momentum conservation) 3 17. Once this foundation is laid, we gradually introduce more chaotic data.  At each step, the network is already in a "good" region of parameter space and is less likely to be derailed by the sharp, complex patterns that chaotic trajectories bring 18 19. This guided progression greatly improves the optimizer's ability to converge to a physically meaningful solution, as evidenced by curriculum learning successes in other non-convex training problems 20. 

Moreover, curriculum learning enhances generalization and robustness 21 5. By first training on high-fidelity, noise-free data, the PINN builds a clean signal of the underlying physics.  Later, when lower-fidelity (slightly noisy) data is added, the network's internal representation already captures the true physics, so it can filter out the numerical noise and avoid overfitting to simulation artifacts 5. In essence, the curriculum acts as a form of regularization: the model must perform well on simple cases before seeing hard cases, which prevents it from simply memorizing chaotic trajectories.  Instead, it must learn general principles that apply across stages.  Given the multi-fidelity nature of our data, transfer learning naturally complements the curriculum.  We do not train a fresh model from scratch on each stage;  instead, we carry over the trained weights from one stage to the next (this combined approach we term Curriculum-Transfer Learning PINN (CTL-PINN)).  Transfer learning leverages the fact that Stage $i+1$ is just Stage $i$ with a bit more difficulty: the solution for the simpler problem is a very good initialization for the harder problem.  This offers multiple advantages:

\begin{itemize}
    \item \textbf{Faster Convergence:} Each time we move to a new stage, the model already "knows" how to handle all easier scenarios, so it can focus on learning the new data's complexities.  This significantly reduces training time on later stages. Empirically, we observe an accelerated convergence at each curriculum step, since the optimizer starts near a good solution rather than from random initialization. 
    \item \textbf{Avoiding Catastrophic Forgetting:} A critical detail in our implementation is replay of earlier data - i.e., when training on Stage i data, we include all samples from stages $0,1,\dots,i-1$ in the training set.  By training on the union of old and new data, the network retains its performance on earlier tasks.  This prevents the common pitfall where a model, upon learning new (harder) cases, forgets what it learned before.  The result is a monotonic expansion of capability: each stage adds to the model's knowledge without degrading prior accuracy. 
    \item \textbf{Implicit Regularization:} Initializing Stage $i$ with weights from Stage $i-1$ biases the training towards solutions that are close to the previous solution in weight space.  This is desirable because the previous solution was physically sensible;  therefore the new solution is constrained not to stray too far into unphysical regions even as it fits more complex data.  In practice, this means the model will seek a solution for the harder stage that still reproduces the simpler-stage behavior for those cases.  Such a bias acts as a regularizer that helps maintain physical consistency and stability across the curriculum. 
\end{itemize}

In summary, the CTL-PINN framework - a staged curriculum with transfer learning and replay - is designed to overcome the training challenges of chaotic dynamics.  It guides the network to first learn conservation laws and stable motion, then gradually introduces chaos in a controlled fashion.  By the final stage, the network has effectively grown up from solving easy Kepler orbits to handling violent three-body interactions, all while maintaining conservation of energy and momentum.  This framework is robust, efficient, and we will demonstrate in the implementation that it is feasible even on modest hardware. 





\subsection{Defining the Curriculum Path}

The curriculum path from stability to chaos is concretely defined by the stage criteria described in Section 2.2 (see Table 1).  To recap briefly, we start exclusively with the high-fidelity hierarchical data in Stage 0, then incrementally loosen the invariant thresholds to include more difficult systems, and eventually incorporate the low-fidelity random data in stages 2 and beyond.  Notably, we chose to introduce the low-fidelity data only after the network has mastered the high-fidelity cases.  An alternative strategy might have been to pre-train on the broader low-fidelity dataset and then fine-tune on the high-fidelity subset.  However, for chaotic dynamics, we deem that approach less effective establishing a correct physics baseline from precise data at the outset is paramount.  If the model were first trained on noisy chaotic data, it might pick up spurious correlations or be misled by numerical errors.  By contrast, our chosen path (HF first, then LF) ensures the PINN has a solid grasp of the true physics (thanks to Stage 0-1 on clean data) and can use that as an internal regularizer when confronted with the lower-precision chaotic cases 22. 

Each stage transition is a point where we expand the training set (adding new simulations) and then continue training (with transferred weights).  We also adjust certain training hyperparameters across stages to reflect the increasing difficulty - for example, the weight on the energy regularization term is ramped up with each stage (details in Section 4.2).  The curriculum is thus not only about data; it also involves a schedule for how strongly we enforce physical constraints as training progresses, ensuring that by the time we reach the most chaotic data, the network is highly constrained to obey invariants.  It's important to highlight that the stage definitions based on invariants are rooted in physical reasoning (bound vs unbound, circular vs eccentric, etc.) but the exact threshold values (e.g. $E<-0.5$ or $\tau_L>5$) are somewhat tunable.  In practice, one could adjust these thresholds or even define more finely graded stages if needed.  The provided five-stage scheme (with a couple of sub-stages) proved sufficient in our experience, striking a good balance between too slow a curriculum (many tiny steps) and too fast (jumps that are too large).  The invariants provide an interpretable way to calibrate these stages.  For example, if the training struggles when Stage 2 data (top 25\% stable random cases) is introduced, one might decide to break that into two smaller introductions (say top 10\% first, then next 15\%).  Conversely, if training is smooth, one could potentially combine some sub-stages.  Thus, while the current stage definitions are mathematically and empirically sound, they are also flexible guidelines. 
\end{comment}
